\documentclass[conference]{IEEEtran}
\usepackage{graphicx}
\usepackage{booktabs}
\usepackage{multirow}
\usepackage{xcolor}
\usepackage{hyperref}
\usepackage{amsmath}
\usepackage{listings}

\newcommand{\openharn}{\textsc{openharn}}
\newcommand{\todo}[1]{\textcolor{red}{\textbf{[TODO: #1]}}}

\title{Inference-Time Recovery of Tool-Use Capability in Low-Bit Quantised Large Language Models}

\author{
\IEEEauthorblockN{Nitesh Sharma}
\IEEEauthorblockA{Independent Researcher}
}

\begin{document}
\maketitle

\begin{abstract}
Small language models (SLMs) are attractive for private, low-cost, and CPU-oriented agent
systems, but tool use fails for reasons that are not captured by ordinary language quality
metrics. A model may select the wrong tool, emit an invalid call, fail to abstain, or produce
one call when a task requires several. We study how much of this failure can be addressed by
an inference-time harness rather than by changing model weights. We describe \openharn, a
small model-agnostic harness with three execution modes: a conventional conversation-based
loop, a grammar-constrained text-call mode, and a structured-action mode with externalized
state and step verification. The harness also provides schema-derived grammars, grounded
filesystem tools, bounded context, text-call recovery, and circuit breakers. On a reproducible
200-example subset of BFCL v4, a 2-bit LFM2 model scores 47.5\% with raw native function
calling. A prompt-tools configuration initially scores 29.5\%, while the later per-request
policy reaches 63.0\% on the primary full-200 run. A subsequent decomposition fix reaches
64.5\% in a later run. On a
separate MiniCPM-V quantization experiment, inference-time settings improve Q4 performance
from 47.5\% to 72.5\% on a 40-example category. These results support a conditional claim:
harnesses can recover mechanical reliability and unlock some latent composition capability when
the presentation matches the model's pretraining distribution, but they cannot supply missing
semantic knowledge or dependent multi-step planning. We provide a failure taxonomy and clearly
separate completed measurements from future evaluation.
\end{abstract}

\section{Introduction}

Small language models are increasingly used on local and resource-constrained hardware.
Their usefulness as agents, however, depends on more than generating plausible text. An agent
must decide whether a tool applies, select the correct tool, emit a valid call, fill its
arguments, execute the call, interpret the result, and stop at the right time. These
requirements expose failure modes that can be hidden by ordinary text-generation scores.

Quantization and model/server compatibility can make the problem worse. Post-training
quantization methods such as GPTQ~\cite{frantar2023gptq} and AWQ~\cite{lin2024awq} reduce
memory and inference cost, but low-bit deployment can expose failures in structured output
behavior. A model may know what operation is needed but fail to emit the format expected by
the server. Conversely, a
grammar can guarantee syntactic validity without making the model choose the right tool or
argument values. Treating all failures as “function calling accuracy” obscures this
boundary.

This paper asks a focused question:
\emph{which tool-use failures can be recovered by an inference-time harness, and which remain
limitations of the model?} We study this question using a compact, CPU-oriented implementation
rather than a new trained model. The harness is designed for small models and exposes the
protocol and execution decisions explicitly enough to inspect them.

Our contributions are:
\begin{enumerate}
\item A systems design for small-model tool use that separates caller-supplied tool schemas
from built-in handlers and supports native calls, prompt-described grammar-constrained calls,
and structured JSON actions.
\item A reliability ladder that combines grounding, schema-derived grammars, text-call
recovery, relevance/abstention handling, bounded context, and circuit breakers. These
mechanisms address different failure classes rather than acting as a universal configuration.
\item An empirical failure taxonomy distinguishing protocol/mechanical failures from
semantic judgment failures, together with measurements on BFCL v4 and a quantization-focused
MiniCPM-V experiment.
\item A conservative account of what the current evidence establishes, plus a controlled
ablation and held-out evaluation protocol needed to test generality beyond the present runs.
\end{enumerate}

\section{Scope and Research Questions}

We study inference-time interventions only: no fine-tuning, distillation, preference
optimization, or model-specific weight changes. The main questions are:

\begin{description}
\item[RQ1:] Can a harness improve valid tool-call generation for a quantized SLM?
\item[RQ2:] Which interventions help single calls, parallel calls, and abstention, and do
those interventions trade off against one another?
\item[RQ3:] Does a compact structured-action loop provide a more reliable alternative when a
model cannot use native function calling?
\item[RQ4:] Which failures remain after syntax and protocol constraints are applied?
\end{description}

The current experiments answer RQ1 and partially answer RQ2 and RQ4. RQ3 is a design
contribution and testable hypothesis: the structured SLM mode is implemented in the artifact,
but it has not been quantitatively compared with the other modes, so this paper makes no
performance claim for it.

\section{Related Work}

The SLM position paper argues that small models can become useful agents when paired with a
good harness~\cite{belcak2025slm}. Our work tests this claim at the tool-call boundary and
qualifies it: a harness can improve form and execution discipline, while semantic judgment
still depends heavily on the model.

ReAct-style agents maintain a conversational trajectory of thoughts, actions, and
observations~\cite{yao2023react}. Toolformer similarly studies training models to decide when
and how to invoke external tools~\cite{schick2023toolformer}. \openharn's default loop follows
this inference-time pattern, while its structured mode externalizes durable state and exposes
only a small valid action set. This design is related to the harness requirements described by
Ranjan and Talluri~\cite{ranjan2026slm}, but our current study does not claim to reproduce their
trained specialist model or their end-to-end results.

BFCL v4 provides AST-based function-calling evaluation without executing arbitrary tools
\cite{patil2025bfcl}. We use its single-turn categories because they isolate call structure,
argument values, parallelism, and abstention. AgentBench provides a broader framework for
evaluating LLM agents across interactive environments~\cite{liu2023agentbench}, while SWE-bench
focuses on issue-resolution performance in real software repositories~\cite{jimenez2024swebench}.
We use BFCL because the present study isolates the tool-call layer rather than end-to-end task
Natural Language Tools~\cite{johnson2025nlt}, tool-necessity prediction~\cite{sun2026probeprefill},
and training-free representation steering~\cite{wang2026asa}, together with grammar-guided
generation~\cite{willard2023guiding}, motivate our prompt-tools, YES/NO selection, and schema
constraint modes. These mechanisms are established individually; our contribution is their
small-model-oriented integration and failure analysis, not the invention of grammar
constraints or tool selection.

\section{Harness Design}

\subsection{Schema and Handler Separation}

Tool schemas are not embedded in the binary. In REPL mode, the caller supplies an OpenAI-style
\texttt{tools} array through \texttt{OPENHARN\_TOOLS\_SCHEMA}; in server mode, schemas arrive
in each request. The harness filters the supplied schemas and derives its prompt and grammar
from them. Execution is separate: a requested function name is dispatched to a built-in Rust
handler. A schema for an unknown name therefore produces an error rather than silently
executing arbitrary code.

The current implementation includes handlers for file reading and writing, anchored editing,
project and system search, shell/Python execution, web fetching, and todo state. Existing
files must be read in the current session before editing or overwriting. This grounding rule
is an execution guard, not a claim that the model understood the whole file.

\subsection{Default Conversational Loop}

The default loop maintains a system prompt, user turns, assistant tool calls, and tool
results. It sends this history to an OpenAI-compatible endpoint, executes accepted calls, and
continues until a text response is produced. The history is trimmed by dropping oldest whole
turns, preserving the system message and recent tool context. Tool results are capped to
prevent one search from consuming the context.

The loop has operational safeguards: per-turn and total call limits, exact-repeat detection,
connection retries, and recovery for structured calls that a server returned as text. Recovery
handles formats such as \texttt{<tool\_call>[...]}; it does not convert arbitrary Markdown
shell fences into tool calls.

\subsection{Grammar-Constrained Text Mode}

Some model/server combinations do not reliably use the native tool API. With
\texttt{OPENHARN\_PROMPT\_TOOLS=1}, tools are described in the system prompt and the native
\texttt{tools} field is omitted. With \texttt{OPENHARN\_STRICT\_TOOLS=1}, a GBNF grammar
restricts output to a schema-valid text call or ordinary text. The grammar is generated from
the caller's schemas, so it constrains function names, argument keys, and value types without
hardcoding a tool inventory.

\textbf{Key systems finding: bounded whitespace.} An unbounded grammar rule allowed a model
to emit a valid first call followed by whitespace until the token limit, leaving no closing
array delimiter. The parser then discarded an otherwise useful call. Bounding this one grammar
rule, together with incomplete-array recovery, produced an approximately 13-point swing in the
measured configuration (about 44\% to 57\%). The intervention is protocol-specific and should
not be treated as a universal model improvement.

\subsection{Structured SLM Mode}

With \texttt{OPENHARN\_SLM=1}, the model does not receive the complete conversational
transcript. It receives a compact JSON observation containing the goal, step budget, valid
actions, discovered files, previous reads, the last result, and any validation feedback. The
action space is deliberately small:
\begin{itemize}
\item \texttt{SEARCH}: search for candidate files;
\item \texttt{READ}: inspect a discovered file;
\item \texttt{ANSWER}: respond after at least one read;
\item \texttt{ESCALATE}: terminate when the task cannot be completed safely.
\end{itemize}

Actions are validated before execution, results are checked afterward, and a failed step is
retried with localized feedback rather than restarting the entire task. This mode tests a
different hypothesis from grammar-constrained tool calling: reducing state and action
complexity may be more useful than teaching a weak model a larger function-call protocol.

\section{Evaluation Protocol}

\subsection{Current Evaluation}

The primary reported benchmark is a fixed 200-example subset of BFCL v4: 40 examples each
from \texttt{simple\_python}, \texttt{multiple}, \texttt{parallel},
\texttt{parallel\_multiple}, and \texttt{irrelevance}. We used this subset because full BFCL
runs were impractical during CPU-only iteration: each condition required local inference and
external AST evaluation, and the parallel categories were especially expensive to repeat.
These are subset results, not BFCL leaderboard scores. The main model is LFM2-8B-A1B-UD-Q2\_K\_XL, a 2-bit quantization with
approximately 1B active MoE parameters. Runs used llama-server on CPU with a 16K context and
12 threads. The BFCL notes report build and temperature details; run-to-run variation is
material, particularly for parallel categories.

A secondary experiment uses MiniCPM-V-4.6 on a 40-example \texttt{parallel\_multiple}
category at Q8 and Q4. It tests whether native template controls can recover quantization
degradation without changing weights.

\subsection{Definition of the Per-Request Policy}

The term ``per-request policy'' refers to a small set of hand-written heuristics implemented
by \texttt{derive\_policy}; it is not a logit router or a learned model. The heuristics inspect
the user text and supplied tool schemas. If the tool list and query indicate one function and
one imperative, the policy prefers native function calling. If the query contains conjunctions
such as ``and'', ``plus'', or ``then'', or contains multiple requested quantities, it treats the
request as multi-call and enables planning. If no schema field matches the query subject, it
uses abstention. In code, \texttt{harness\_decompose} estimates the number of requested
operations and selects a generation path:
\begin{itemize}
\item For a likely single-call request, it prefers native function calling and does not force
prompt-tools or a call-count hint.
\item For a likely multi-call request, it tries the model's native template with a think-then-
call phase, then a plan-first prompt that enumerates operations before emitting a call array,
and finally strict prompt-tools with an optional count hint.
\item For an irrelevant request, or when the decomposition finds no matching tool, it uses the
relevance/abstention path and can return \texttt{NO\_TOOL} without another model call.
\end{itemize}

Run D is this policy plus two proxy fixes: flattening BFCL's double-nested message structure
before sending it to llama-server, and retrying with flattened prompt-tools plus strict grammar
when native function calling returns no call. Run D therefore corresponds to the
``per-request policy'' row in Table~\ref{tab:bfcl}; it does not mean that every request was
forced through the same H2, deduplication, or prompt-tools pipeline.

Run E adds the decomposition correction and fast path that were not present in Run D. The
original decomposition code expected OpenAI-wrapped tool objects, whereas BFCL supplied flat
objects; consequently it silently estimated zero matching tools. The fix accepts the flat
schema, makes the no-match path abstain immediately, and simplifies the relevance-gate prompt
with a ``when in doubt, say no'' instruction. This is the decomposition fix behind the 64.5\%
run. Incomplete-array recovery and bounded-whitespace grammar are shared parser/grammar
components, but they are not what distinguishes Run E from Run D.

\subsection{Fairness and Reproducibility Requirements}

Configurations were developed iteratively on the fixed subset; reported numbers are therefore
directional rather than unbiased estimates. A rigorous methodology would freeze prompts, schemas,
decoding parameters, and model/server versions before scoring held-out examples, and would report
repeated runs with confidence intervals. It should also compare each intervention independently
and in combination, using metrics such as latency, generated tokens, invalid-call rate,
hallucinated-result rate, and final task success.

\section{Results}

\subsection{BFCL Results}

Table~\ref{tab:bfcl} summarizes the results documented in the BFCL research notes. The raw
native-function-calling result is the baseline. Prompt-tools plus strict grammar is not
automatically better: the initial configuration scored lower because the model's native
training and the text protocol competed. The later per-request policy, including message
flattening and native-empty fallback, produced 63.0\% on the full 200-example run. A
subsequent decomposition fix produced 64.5\%; we use 63.0\% as the primary result because it
is the stable run discussed in the research notes.

\begin{table}[t]
\centering
\caption{BFCL v4 results on the fixed 200-example subset. Values are percentages. Run D is the per-request policy plus BFCL message flattening and native-empty fallback; Run E reached 64.5\% after an additional decomposition fix.}
\label{tab:bfcl}
\resizebox{\columnwidth}{!}{%
\begin{tabular}{lrrrrrr}
\toprule
Configuration & Simple & Multiple & Parallel & Parallel-M & Irrel. & Overall\\
\midrule
Raw native FC & 80.0 & 82.5 & 0.0 & 0.0 & 75.0 & 47.5\\
Prompt + strict & 35.0 & 10.0 & 5.0 & 5.0 & 92.5 & 29.5\\
+ abstain + gate & 60.0 & 45.0 & 2.5 & 25.0 & 87.5 & 44.0\\
Per-request policy (Run D) & 75.0 & 77.5 & 52.5 & 42.5 & 67.5 & 63.0\\
\bottomrule
\end{tabular}%
}
\end{table}

The main pattern is more informative than the aggregate score. Native calling is competitive
for ordinary single calls but produces no correct parallel calls for this model and setup. The
native-template and plan-first paths improve some multi-call cases by giving the model a planning
phase before enforcing an array, while the relevance gate improves abstention. The same strict prompt-tools configuration can damage models whose
native function-calling path is healthy; measurements in the project notes show strong native
models falling to zero on the tested parallel-multiple setup under the incompatible strict
configuration. Thus the result supports per-model or per-request adaptation, not one
universal flag set.

\subsection{Quantization Experiment}

As a cross-model check, we tested whether inference controls can recover quantization
degradation on a different model family, MiniCPM-V. Table~\ref{tab:quant} reports the separate
experiment. The Q4 model's score rises
from 47.5\% with automatic native tool choice to 72.5\% when tool choice is required and
thinking is disabled. The result is encouraging but narrow: it is one model family, one
category, and a small evaluation slice. It demonstrates that inference controls can matter
substantially; it does not establish that quantization primarily damages format rather than
other capabilities.

\begin{table}[h]
\centering
\caption{MiniCPM-V-4.6 quantization experiment on 40 parallel-multiple examples.}
\label{tab:quant}
\begin{tabular}{lrr}
\toprule
Configuration & Accuracy & Repeated values\\
\midrule
Q8 native, automatic & 82.5\% & 77.5/85/85\\
Q4 native, automatic & 47.5\% & 45/50/47.5\\
Q4 required + no-think & 72.5\% & 72.5/72.5/72.5\\
\bottomrule
\end{tabular}
\end{table}

\subsection{What the Current Results Do Not Establish}

The present experiments do not establish a universal improvement, a causal ranking of every
harness component, or superiority over trained specialist agents. The 200-example BFCL slice
was selected for practical feasibility, and some configurations were tuned against related
behavioral cases. The results should therefore be read as evidence for specific mechanisms
and failure modes, not as a leaderboard claim.

\section{Failure Taxonomy}

We classify failures by the layer at which they occur:

\begin{itemize}
\item \textbf{Protocol failure.} The model emits a format the server cannot parse, or the
grammar/parser loses a call through truncation or runaway whitespace. \textbf{The bounded-
whitespace grammar fix is a concrete systems result: replacing an unbounded whitespace rule
with a one-space bound prevented runaway output and produced an approximately 13-point swing
between the affected BFCL configurations.} Recovery can help.
\item \textbf{Selection or abstention failure.} The model calls a tool when none applies, or a
relevance gate says no when a tool is needed. A gate can improve abstention but introduces
false negatives.
\item \textbf{Argument failure.} The model selects a plausible function but supplies a wrong
value, type, enum, or nested argument. A grammar can enforce types, not facts.
\item \textbf{Planning failure.} The model selects the wrong number of calls, under-decomposes a
parallel task, duplicates calls, or stops at the wrong time. This remains primarily model
judgment, although prompts and structured state can change the interface.
\end{itemize}

This taxonomy prevents a common overclaim. Valid JSON is not equivalent to correct tool use,
and a higher tool-call score does not necessarily imply better end-to-end coding behavior.
The appropriate metrics should report syntax validity, tool selection, argument accuracy,
call-count accuracy, abstention, execution success, and final task success separately.



\section{Discussion}

The evidence supports a qualified systems conclusion. Harnesses can recover mechanical
reliability: they can constrain syntax, prevent unsafe edits, express multiple calls, repair
some truncation, and supply an explicit abstention mechanism. They cannot reliably provide
missing semantic knowledge, correct argument values, tool competence, or robust decomposition.

The best configuration is therefore conditional. A model with healthy native function calling
should not automatically be forced through prompt-tools and strict grammar. A weak model that
cannot emit native calls may benefit from a narrow schema, prompt descriptions, and grammar
constraints. A model that can emit compact JSON but not tool calls may be a candidate for the
structured SLM mode. The practical design principle is a reliability ladder: add constraints
only when measurements show that the corresponding failure exists.

The CPU results also expose a deployment tradeoff. Reasoning tokens can dominate wall-clock
latency, especially on local hardware. Although the project notes suggest that disabling
thinking can improve throughput, exact latency values were not retained for the MiniCPM-V
tabled runs, so we do not claim a measured speedup here. Future runs should report the
speed/accuracy tradeoff directly.

\subsection{Limitations and Threats to Validity}

\begin{itemize}
\item The primary BFCL result is a fixed 200-example subset, not the full leaderboard.
\item The main BFCL model is one LFM2 quantization and the secondary experiment is one
MiniCPM-V family; broad transfer is untested.
\item Configurations were developed iteratively on the fixed subset, so reported numbers are
directional rather than unbiased estimates; a full held-out multi-model ablation is future work.
\item The measurements depend on llama-server chat templates, grammar behavior, and build
versions. Another server may produce different results. Guided-generation implementations
also differ in their supported grammars and decoding semantics~\cite{willard2023guiding}.
\item BFCL single-turn AST accuracy is not end-to-end coding success and does not evaluate
filesystem safety, multi-turn state, or final answer quality.
\item The structured SLM implementation is quantitative evidence only after a direct controlled
comparison; its current inclusion is a design contribution and testable hypothesis.
\item The artifact contains powerful local handlers such as shell and Python execution. The
research evaluation assumes a trusted local environment and does not claim sandbox security.
\end{itemize}

\section{Conclusion}

We presented a small-model-oriented tool-use harness and a framework for analyzing where it
helps. Existing measurements show that schema-derived grammar, abstention handling, bounded
output, and inference controls can substantially change tool-call reliability without training.
They also show that a prompt-tools configuration can hurt a model whose native interface is
already working, and that semantic judgment remains a bottleneck after syntax is fixed.

The defensible conclusion is not that a harness is universally more important than a model.
It is that harness design is a major, measurable part of small-model agent capability. A good
harness can recover form and execution discipline; the model still determines much of the
judgment. The next step is the controlled, held-out, multi-model ablation described above.

\section*{Acknowledgments}

We thank the BFCL team and the Gorilla project for the BFCL v4 datasets and evaluation code,
and the llama.cpp community for the OpenAI-compatible inference interface used in the
experiments. The reported runs were conducted on consumer CPU-oriented hardware.

\bibliographystyle{IEEEtran}
\bibliography{references}

\end{document}
